\documentclass[12pt]{article}
\usepackage[utf8]{inputenc}
\usepackage[T1]{fontenc}
\usepackage[margin=1in]{geometry}
\usepackage{setspace}
\usepackage{xcolor}
\usepackage[colorlinks=true, urlcolor=blue!60!black, linkcolor=blue!60!black]{hyperref}
\setlength{\parindent}{0em}
\setlength{\parskip}{0.6em}
\renewcommand{\baselinestretch}{1.0}
\begin{document}
\begin{center}
{\LARGE \textit{Cohen v.\ California}, 403 U.S.\ 15 (1971)\par}
\vspace{0.5em}
{\large Case Brief --- Two-Page Edition\par}
\end{center}

\section{Memory Jogger}
California could not punish a jacket reading ``Fuck the Draft'' in a courthouse corridor: not obscenity, not fighting words, and viewers could avert their eyes (\href{https://supreme.justia.com/cases/federal/us/403/15/}{Justia}).

\section{Detailed Case Facts}
In 1968, nineteen-year-old Paul Robert Cohen wore a jacket bearing ``Fuck the Draft'' in a corridor of the Los Angeles County Courthouse, where women and children were present, to show the depth of his feelings against the Vietnam War. He made no noise, threatened no one, and provoked no violence.

\section{Procedural History}
Cohen was convicted of disturbing the peace by ``offensive conduct'' and sentenced to 30 days. The California Court of Appeal affirmed, reasoning that the jacket might provoke violence; the California Supreme Court denied review. The U.S.\ Supreme Court reversed.

\section{Judicial Votes}
Reversed 5--4. Harlan wrote for Douglas, Brennan, Stewart, and Marshall; Blackmun dissented, joined by Burger and Black, and by White in part.

\section{Holding}
Absent a more particularized and compelling reason, a State may not criminalize the public display of a single four-letter expletive. Profanity conveying a political message is protected, and passersby in a public corridor are not a captive audience.

\section{Analysis of Opinions}
\textbf{Majority:} the conviction rested on words alone, not conduct; the statute was no courthouse-decorum rule; the slogan was not erotic, not directed at anyone present, and could be avoided ``simply by averting their eyes.'' Officials cannot draw principled lines among offensive words, emotive force deserves protection, and word-censorship risks suppressing ideas. \textbf{Blackmun:} the ``absurd and immature antic'' was ``mainly conduct, and little speech,'' and the case should be remanded under \textit{In re Bushman}.

\section{Examples: Future Applications}
\textbf{Same-Side:} a profane protest sign outside city hall. \textbf{Opposite-Side:} a face-to-face epithet daring a rival to fight (\textit{Chaplinsky}). \textbf{Fence-Sitter:} a narrow ban on profane apparel inside a courtroom in session.

\section{Critique}
Communitarian critics say ``avert your eyes'' undervalues public decorum in civic spaces; others note the Court later drew lines around profanity in broadcasting (\textit{Pacifica}) and schools (\textit{Fraser}).

\section{Key Quotations}
\begin{quote}
``For, while the particular four-letter word being litigated here is perhaps more distasteful than most others of its genre, it is nevertheless often true that one man's vulgarity is another's lyric. Indeed, we think it is largely because governmental officials cannot make principled distinctions in this area that the Constitution leaves matters of taste and style so largely to the individual.''
\end{quote}
\noindent --- Justice Harlan, for the Court, 403 U.S.\ at 25.

\end{document}
